% Part: incompleteness
% Chapter: arithmetization-syntax
% Section: proofs-in-ax

\documentclass[../../../include/open-logic-section]{subfiles}

\begin{document}

\olfileid[pt]{inc}{art}{pax}
\olsection{\usetoken{P}{derivation} axiomáticas}

\begin{explain}
Para aritmetizar !!{derivation}s axiomáticas, precisamos representar
!!{derivation}s como números. Como !!{derivation}s são simplesmente sequências de
!!{formula}s, a abordagem evidente é codificar cada !!{derivation}
como o código da sequência dos códigos das !!{formula}s nela contidas.
\end{explain}

\begin{defn}
Se $\delta$ é !!a{derivation} axiomática que consiste nas !!{formula}s
$!A_1$, \dots,~$!A_n$, então $\Gn{\delta}$ é
\[
\tuple{\Gn{!A_1}, \dots, \Gn{!A_n}}.
\]
\end{defn}

\begin{ex}
  Considere a !!{derivation} muito simples:
  \begin{derivation}
  1. & $!B \lif (!B \lor !A)$ \\
  2. & $(!B \lif (!B \lor !A)) \lif (!A  \lif (!B \lif (!B \lor !A)))$\\
  3. & $!A  \lif (!B \lif (!B \lor !A))$
  \end{derivation}
  O número de G\"odel dessa !!{derivation} seria
  \begin{align*}
  \openTuple\, 
    & \Gn{!B \lif (!B \lor !A)}, \\
    &\Gn{(!B \lif (!B \lor !A)) \lif (!A  \lif (!B \lif (!B \lor !A)))},\\
    & \Gn{!A  \lif (!B \lif (!B \lor !A))} \,\closeTuple.
  \end{align*}
\end{ex}

\begin{explain}
Depois de fixar uma representação de !!{derivation}s, também precisamos
mostrar que podemos manipular tais !!{derivation}s de maneira recursiva
primitiva e expressar assim suas propriedades e relações essenciais.
Algumas operações são simples: p.ex., dado um número de
G\"odel~$d$ de !!a{derivation}, $(d)_{\len{d}-1}$ nos dá o número de
G\"odel de sua !!{formula} final. Algumas são muito mais difíceis.
Esboçaremos ao menos como fazer isso. O objetivo é mostrar que a
relação ``$\delta$ é !!a{derivation} de~$!A$ a partir de~$\Gamma$'' é
recursiva primitiva nos números de G\"odel de $\delta$ e~$!A$.
\end{explain}

\begin{prop}
\ollabel{prop:followsby}
As seguintes relações são recursivas primitivas:
\begin{enumerate}
\item $!A$ é um axioma.
\item A $i$-ésima linha de $\delta$ é justificada por modus ponens.
\item A $i$-ésima linha de $\delta$ é justificada por \QR.
\item $\delta$ é uma !!{derivation} correta.
\end{enumerate}
\end{prop}

\begin{proof}
Temos de mostrar que as relações correspondentes entre números de
G\"odel de !!{formula}s e números de G\"odel de !!{derivation}s são
recursivas primitivas.
\begin{enumerate}
\item Temos uma lista dada de esquemas de axiomas, e $!A$ é um axioma
  se tem a forma dada por um desses esquemas. Como a lista de esquemas
  é finita, basta mostrar que podemos testar recursivamente de maneira
  primitiva, para cada esquema de axiomas, se $!A$ tem essa forma. Por
  exemplo, considere o esquema de axiomas
  \[
  !B \lif (!C \lif !B).
  \]
  $!A$ é uma instância desse esquema de axiomas se existem !!{formula}s
  $!B$ e $!C$ tais que obtemos $!A$ ao concatenar `$($' com $!B$, com
  `$\lif$', com `$($', com $!C$, com `$\lif$', com $!B$ e com~`$))$'.
  Podemos testar a propriedade correspondente do número de G\"odel
  $n$ de~$!A$, pois a concatenação de sequências é recursiva primitiva
  e os números de G\"odel de $!B$ e $!C$ devem ser menores que o
  número de G\"odel de~$!A$: quando a relação vale, tanto $!B$ quanto
  $!C$ são sub-!!{formula}s de~$!A$. Portanto, podemos definir:
  \begin{multline*}
  \fn{IsAx}_{!B \lif (!C \lif !B)}(n) \defiff \bexists{b< n}{
    \bexists{c < n}{(\fn{Sent}(b) \land \fn{Sent}(c) \land {}}}\\ n =
  \Gn{(} \concat b \concat \Gn{\lif} \concat \Gn{(} \concat c \concat
  \Gn{\lif} \concat b \concat \Gn{))}).
  \end{multline*}
  Se tivermos uma definição assim para cada esquema de axiomas, a
  disjunção delas define a propriedade $\fn{IsAx}(n)$, ``$n$~é o
  número de G\"odel de um axioma''.
\item A $i$-ésima linha de $\delta$ é justificada por modus ponens sse
  existem linhas $j$ e $k < i$ nas quais a !!{sentence} da linha~$j$
  é alguma fórmula $!A$, a sentença da linha~$k$ é $!A \lif
  !B$, e a sentença da linha~$i$ é~$!B$.
  \begin{multline*}
    \fn{MP}(d, i) \defiff \bexists{j < i}{\bexists{k < i}{}}\\
    (d)_k = \Gn{(} \concat (d)_j
      \concat \Gn{\lif} \concat (d)_i \concat \Gn{)}
  \end{multline*}
  Como a quantificação limitada, a concatenação e $=$ são recursivas
  primitivas, isso define uma relação recursiva primitiva.
\item Uma linha de $\delta$ é justificada por \QR{} se tem a forma
  $!B \lif \lforall[x][!A(x)]$, uma linha anterior é $!B \lif !A(c)$
  para alguma !!{constant}~$c$, e $c$ não ocorre em~$!B$. Esse é o
  caso sse
  \begin{enumerate}
  \item existe !!a{sentence}~$!B$ e
  \item !!a{formula}~$!A(x)$ com uma única variável~$x$ livre, tais que
  \item a linha $i$ contém $!B \lif \lforall[x][!A(x)]$,
  \item alguma linha $j < i$ contém $!B \lif \Subst{!A}{c}{x}$ para
    uma constante~$c$,
  \item que não ocorre em~$!B$.
  \end{enumerate}
  Tudo isso pode ser testado recursivamente de maneira primitiva, pois
  os números de G\"odel de $!B$, $!A(x)$ e $x$ são menores que o número
  de G\"odel da fórmula na linha~$i$, e o de $a$ é menor que o
  número de G\"odel da fórmula na linha~$j$:
  \begin{multline*}
    \fn{QR}_1(d, i) \defiff \bexists{b < (d)_i}{\bexists{x <
        (d)_i}{\bexists{a < (d)_i}{\bexists{c < (d)_j}{(}}}}
    \\ \fn{Var}(x) \land \fn{Const}(c) \land {} \\
    (d)_i = \Gn{(} \concat
    b \concat \Gn{\lif} \concat \Gn{\lforall} \concat x \concat a
    \concat \Gn{)} \land {}\\
    (d)_j = \Gn{(} \concat b \concat
    \Gn{\lif} \concat \fn{Subst}(a,c,x) \concat \Gn{)} \land {}\\ \fn{Sent}(b)
    \land \fn{Sent}(\fn{Subst}(a,c,x)) \land {} \bforall{k <
      \len{b}}{(b)_k \neq (c)_0})
  \end{multline*}
  Aqui supomos que $c$ e $x$ são os números de G\"odel da variável e
  da constante consideradas como termos (O.s.\@, não seus códigos de
  símbolos). Testamos se $x$ é a única variável livre de~$!A(x)$
  testando se $\Subst{!A(x)}{c}{x}$ é !!a{sentence}, e asseguramos que
  $c$ não ocorre em $!B$ exigindo que todo símbolo de~$!B$ seja
  diferente de~$c$.

  Deixamos a outra versão de \QR{} como exercício.
\item $d$ é o número de G\"odel de uma !!{derivation} correta sse cada
  linha nela é um axioma ou é justificada por modus ponens ou~\QR.
  Portanto:
  \[
  \fn{Deriv}(d) \defiff \bforall{i < \len{d}}{(\fn{IsAx}((d)_i) \lor
    \fn{MP}(d,i) \lor \fn{QR}(d, i))}
  \]
\end{enumerate}
\end{proof}

\begin{prob}
Defina as seguintes relações como em
\olref[inc][art][pax]{prop:followsby}:
\begin{enumerate}
\item $\fn{IsAx}_{!A \lif (!B \lif (!A \land !B))}(n)$,
\item $\fn{IsAx}_{\lforall[x][!A(x)] \lif !A(t)}(n)$,
\item $\fn{QR}_{2}(d, i)$ (para a outra versão de \QR).
\end{enumerate}
\end{prob}

\begin{prop}
Suponha que $\Gamma$ seja um conjunto recursivo primitivo de
!!{sentence}s. Então a relação $\Prf[\Gamma](x, y)$ que expressa
``$x$ é o código de !!a{derivation}~$\delta$ de $!A$ a partir de~$\Gamma$
e $y$ é o número de G\"odel de~$!A$'' é recursiva primitiva.
\end{prop}

\begin{proof}
Suponha que ``$y \in \Gamma$'' seja dado pelo predicado recursivo
primitivo~$R_\Gamma(y)$. Temos de mostrar que a relação
$\Prf[\Gamma](x, y)$ é recursiva primitiva, em que $\Prf[\Gamma](x,
y)$ vale sse $y$ é o número de G\"odel de !!a{sentence}~$!A$ e
$x$~é o código de !!a{derivation} de~$!A$ a partir de $\Gamma$.

Pela proposição anterior, a propriedade $\fn{Deriv}(x)$, que vale sse
$x$ é o código de uma !!{derivation} correta~$\delta$, é recursiva
primitiva. Contudo, essa definição não levou em conta o conjunto
$\Gamma$ como uma maneira adicional de justificar linhas na
!!{derivation}. Nosso teste recursivo primitivo de se uma linha é
justificada por \QR{} tampouco levou em conta a exigência de que a
constante~$c$ não possa ocorrer em~$\Gamma$. É possível modificar
nossa definição para que leve $\Gamma$ diretamente em conta, mas é
mais fácil usar $\fn{Deriv}$ e o teorema da dedução. $\Gamma \Proves
!A$ sse existe alguma lista finita de !!{sentence}s $!B_1$, \dots,
$!B_n \in \Gamma$ tal que $\{!B_1, \dots, !B_n\} \Proves !A$. E, pelo
teorema da dedução, esse é o caso se $\Proves (!B_1 \lif (!B_2 \lif
\cdots (!B_n \lif !A)\cdots))$. Pode-se testar recursivamente de
maneira primitiva se !!a{sentence} com número de G\"odel~$z$ tem essa
forma. Assim, em vez de considerar $x$ como o número de G\"odel de
!!a{derivation} da !!{sentence} com número de G\"odel~$y$ \emph{a
partir de $\Gamma$}, consideramos $x$ como o número de G\"odel de
!!a{derivation} de uma condicional aninhada da forma acima a partir
de~$\emptyset$.

Primeiro, se temos uma sequência de !!{sentence}s, podemos formar
recursivamente de maneira primitiva a condicional que tem todas essas
sentenças como antecedentes e uma dada !!{sentence} como consequente:
\begin{align*}
  \fn{hCond}(s, y, 0) & = y \\
  \fn{hCond}(s, y, n+1) & = \Gn{(} \concat (s)_{n} \concat \Gn{\lif}
  \concat \fn{Cond}(s, y, n) \concat \Gn{)}\\
  \fn{Cond}(s, y) & = \fn{hCond}(s, y, \len{s})\\
  \intertext{Assim, podemos definir $\Prf[\Gamma](x, y)$ por}
  \Prf[\Gamma](x, y) & \defiff \bexists{s < \fn{sequenceBound}(x,x)}{(} \\
  &\qquad (x)_{\len{x}-1} = \fn{Cond}(s,y) \land {} \\
  &\qquad\bforall{i<\len{s}}{(s)_i \in \Gamma} \land {} \\
  &\qquad\fn{Deriv}(x)).
\end{align*}
O limite de~$s$ é dado pela observação de que cada $(s)_i$ é o número
de G\"odel de uma sub-!!{formula} da última linha da !!{derivation},
O.s.\@, é menor que $(x)_{\len{x}-1}$. O número de antecedentes $!B
\in \Gamma$, O.s.\@, o comprimento de~$s$, é menor que o comprimento
da última linha de~$x$.
\end{proof}

\end{document}
